\documentclass[11pt]{article}
\usepackage[a4paper,margin=25mm]{geometry}
\usepackage{amsmath,amssymb,amsthm,hyperref}
\newtheorem{theorem}{Theorem}
\newcommand{\Z}{\mathbb Z}\newcommand{\F}{\mathbb F}\newcommand{\Q}{\mathbb Q}
\title{E.214: an original-family class-three survivor with empty class four}
\author{Recovery transcription}\date{18 November 2026}
\begin{document}\maketitle
\paragraph{Recovery provenance.}
This is a complete mathematical transcription of the reviewed E.214 proof, reconstructed after loss of the working filesystem. It is not a claim of byte identity with the lost canonical TeX. The original proof received complete positive independent reviews from root and ghost\_audit. Its original checker was independently executed by root; the recorded JSON SHA256 was
\begin{center}\texttt{3c6351684e6d3a3b1ad939d51308aae587ff1aef2e5ad3dae08b6cffb67d3f84}.\end{center}
This transcription neither recreates that execution nor substitutes a new checker for the original one.

\section*{Statement and precise scope}
Use the original A.44 factorial construction, its legal root choices, and the minimal five-adic shells of E.199. Thus
\[
 n_j=(j+4)!,\qquad c_j=v_5(n_j),\qquad
 e_j=v_5(n_{j+1})=c_j+v_5(j+5),
\]
and at the new A- and B-support primes one has
\[
 v_5(q_j-1)=v_5(r_j-1)=e_j.
\]
The initial exponent is $n_0=24$. The older rational valuations at $3$ and $2$ are respectively $(1,0)$ and $(0,1)$. The completed pair is denoted $(A,B)$.

\begin{theorem}[E.214]
There are outputs of the original A.44 recursion satisfying all of its original arithmetic conditions for which the integral pro-five class-two section is matched to the ordinary point class at every off-diagonal support prime and the global integral class-three fibre is nonempty, but the global integral class-four fibre is empty. One can impose the two obstructions at $q_{10}$ and $q_{15}$. Consequently class four is the first empty integral pro-five nilpotent level for these outputs.

This is a selected subfamily statement. It neither excludes every output of A.44 nor proves anything about survival of a fixed pair through all nilpotent levels or all coefficient primes.
\end{theorem}

\begin{proof}
\subsection*{1. Established inputs and an actual finite representation}
We use E.213 to choose a legal first pair with consistent first five-adic gate, and E.211 to continue it inside the unchanged A.44 recursion while matching the integral class-two section at every support prime. E.209 gives
\[
 H^2(G_{\Q,\{5,\infty\}},\Z_5(w))=0\quad(w\ge2),
\]
and E.210 gives
\[
 H^1(G_\Q,\Z_5(2))=0,\qquad
 H^1(G_\Q,\Z_5(3))=\Z_5\epsilon.
\]
The class-two section is therefore unique over a fixed integral abelian pair, up to the prescribed geometric conjugacy. The matched-inertia comparison of B.25 then supplies a global integral class-three lift. These inputs assert existence of an actual Galois section; no formal Lie action is substituted for it.

Fix one legal prefix through $q_6,r_6$, and choose one E.211 continuation and one actual integral class-three lift $t_*$. Let $\bar t_*$ be its reduction in the class-three exponent-five quotient. Since $3<5$, this quotient is the Lazard group of the free class-three Lie algebra over $\F_5$. Write
\[
 C=[X,Y],\qquad D=[C,X],\qquad E=[C,Y].
\]
It has order $5^5$. The actual arithmetic action is retained. In the tangential normalization its degree-two path coefficient is
\[
 \eta(g)=\frac{\chi(g)^2-1}{24}.
\]
Thus the scalar part $\lambda=\bar\chi(g)\in\F_5^\times$ acts on the degree-one generators by
\[
 X\longmapsto\lambda X,\qquad
 Y\longmapsto\lambda Y-\lambda\bar\eta(g)E,
\]
and on $C,D,E$ by their graded weights $2,3,3$. This is the reduction of the actual based action; $24$ is a five-adic unit.

For an element with zero $Y$ coordinate and nonzero $X$ coordinate put
\[
 L=\alpha X+cC+dD+eE,\qquad
 \mathcal K(L)=e-\frac{c^2}{2\alpha}.
 \tag{1}
\]
This quantity is invariant under geometric conjugacy. For example conjugation by $\exp(-uY)$ changes
\[
 c\mapsto c+\alpha u,\qquad
 e\mapsto e+cu+\frac{\alpha u^2}{2},
\]
which leaves (1) unchanged; conjugations in the other Hall directions do not change (1). Arithmetic conjugation multiplies $\mathcal K$ by $\lambda^3$. The reduction of $\epsilon$ transforms by the same factor. Hence the ratio $\mathcal K/\epsilon(F)$, whenever defined, is invariant under the conjugacy ambiguity in a Chebotarev prescription.

\subsection*{2. The common finite class-three torsor after the prefix}
Consider only continuations retaining the chosen prefix through $q_6,r_6$ and choosing principal-unit ratios at every already corrected support place. This is allowed by E.211. Every later exponent is divisible by $25$; indeed the first later exponent is $n_7=11!$. More importantly, both old A- and B-coordinate primitive ratios are retained as principal units. This is stronger than merely requiring equality of their Kummer classes modulo $25$.

Choose common literal abelian cocycles modulo $25$. The finite class-two differences have zero inertia away from $5$: at later support the exponents are divisible by $25$, at the fixed old support the principal-unit comparison gives the same inertia, and at zero support the integral tame relation gives zero inertia. The fixed-ramification calculation
\[
 H^1(G_{\Q,\{5,\infty\}},\Z/25(2))=0
\]
therefore permits a global central coboundary normalization making the class-two cochains identical modulo $25$.

The reduction modulo $5$ of the class-three difference is unramified away from $5$. We spell out the delicate first-support point. At $q_0$ or $r_0$ the class-two aligning central coordinate is a divided quantity with denominator of five-adic valuation one. Literal equality of the class-two Frobenius cochain modulo $25$, together with the stronger $25$-divisibility of the old weight-one primitive changes, makes that aligning coordinate change by a multiple of $5$. Thus the point inertia and every integral class-three lift have the same reduction there. At $2$ and $3$ the relevant weight-two and weight-three denominators are units. At every later support the class-two section is matched to the point; every integral class-three lift differs from a point lift by an unramified weight-three cocycle. The point inertia is a fifth power and is trivial in the exponent-five quotient. Zero-support inertia again vanishes by the integral tame relation.

Consequently every actual integral class-three lift of every retained continuation has reduction
\[
 \bar t_*+\bar\epsilon\,(\lambda_DD+\lambda_EE),
 \qquad (\lambda_D,\lambda_E)\in\F_5^2.
 \tag{2}
\]
This conclusion concerns actual integral lifts, not all possible formal defining systems. Conversely we only need the previously established nonemptiness of the integral class-three fibre.

\subsection*{3. Full geometric image and the exact radical intersection}
Retain the finite field cut out by the joint actual representation consisting of $\bar t_*$, the actual class-three action and cusp-path data, and $\bar\epsilon$. Over $\Q(\mu_5)$ its geometric part is
\[
 P\times\F_5(3),
\]
where $P$ is the free class-three exponent-five group. To see fullness, inertia at $3$ and at $2$ acts trivially on the based arithmetic data and on $\epsilon$ and gives a Frattini basis of $P$: the abelian valuations are $(1,0)$ and $(0,1)$. Thus the kernel of the arithmetic projection maps onto $P$. The nonzero class $\bar\epsilon$ then supplies the independent last factor. The full group is the corresponding semidirect product with the actual cyclotomic $C_4$ action. Its maximal abelian subfield is $\Q(\mu_5)$, since the weight-one geometric generators and the weight-three $\epsilon$ direction lie in the derived subgroup.

At a new A-prime, let $F$ be the full cyclotomic CRT field in the original recursion, and let
\[
 R=F\bigl(B_{\mathrm{current}}^{1/n_{j+1}}\bigr)
\]
be its required complementary radical field (the equivalent original reciprocal convention gives the same argument). All root and old-local congruence prescriptions are retained. Over $F$ the intersection with the actual finite five-field is cyclic of exponent at most five. Any cyclic quotient of $P$ factors through its abelianization, with cyclotomic weight one; the $\epsilon$ line has weight three and cannot be a rational Kummer radical line. The two rational Kummer lines are independent over the full cyclotomic field: the private valuations at $3$ and $2$ are prime to five, and the relevant local cyclotomic ramification indices there are prime to five. The unique possible intersection is therefore the complementary fifth-root line
\[
 F\bigl(B_0^{1/5}\bigr).
 \tag{3}
\]
Here $B_0$ is the retained complementary Kummer class; later factorial powers do not change this line. This identifies the whole cyclic intersection, rather than treating the class-three field as abelian.

Prescribing zero $Y$ coordinate fixes the marked root in (3). It is consequently compatible with the original fixed-root radical automorphism. The remaining original CRT and old-root data lie in $F$ or in the prime-to-five part of the radical field and are untouched. This proves that the following are genuine joint Frobenius prescriptions, not merely elements of an abstract Lie group.

\subsection*{4. Two legal, incompatible transverse prescriptions}
At the A-prime $q_{10}$ prescribe
\[
 \bar t_*(F_{q_{10}})=\exp X,\qquad
 \bar\epsilon(F_{q_{10}})=1,
\]
and at $q_{15}$ prescribe
\[
 \bar t_*(F_{q_{15}})=\exp(X+E),\qquad
 \bar\epsilon(F_{q_{15}})=1.
 \tag{4}
\]
Both have zero complementary coordinate, nonzero A-coordinate, and the required actual arithmetic component. Chebotarev in the compositum just described realizes them together with all original A.44 constraints, arbitrary size requirements, and the minimal-five shell. The E.211 delayed root correction remains available for every such choice. Its previously corrected roots stay principal, so the common-torsor argument above remains valid.

For a local A-point, let $u$ be the weight-one complementary primitive and let
\[
 R_A=\frac{h(F_q)-q\alpha u}{q^2-1}
\]
be the actual integral class-two aligning coordinate. In the retained Lie convention the aligned point's class-three coordinates are
\[
 c_{\mathrm{pt}}=\alpha u,\qquad
 d_{\mathrm{pt}}=-\alpha R_A,\qquad
 e_{\mathrm{pt}}=\frac{\alpha u^2}{2}.
 \tag{5}
\]
Thus its transverse quantity (1) is zero. In particular the divided coordinate $R_A$ has not been discarded: it affects the longitudinal component, while (1) is the exact transverse invariant.

For a general integral class-three lift (2), equations (4)--(5) give transverse discrepancies
\[
 \lambda_E\quad\text{at }q_{10},\qquad
 1+\lambda_E\quad\text{at }q_{15}.
 \tag{6}
\]
In this proof $Z$ denotes \emph{section minus aligned point}; it has the opposite sign from B.26's required-correction variable. With this convention its transverse coordinate is $\mathcal K$, not $-\mathcal K$. The zero/nonzero conclusion is independent of that convention.

\subsection*{5. Why a nonzero discrepancy blocks every class-four lift}
Work integrally in the class-four Lazard group at coefficient five; all required denominators are units. After alignment with the actual point, write the local inertia and Frobenius in the form
\[
 T=\exp(mX+V),\qquad
 F=\exp Z\,F_{\mathrm{pt}}\exp W,
\]
where $Z\in L_3$ is the unramified class-three discrepancy and $V,W\in L_4$. The actual tame relation is $FTF^{-1}=T^q$. Its weight-four part is
\[
 (q^4-q)V+qm[Z,X]=0,
\]
or, since $q$ is a five-adic unit,
\[
 (q^3-1)V+m[Z,X]=0.
 \tag{7}
\]
The Frobenius correction $W$ is central and drops out. The map
\[
 \operatorname{ad}_X:L_3\longrightarrow L_4
\]
is a primitive injection: its two Hall basis vectors map to independent Hall vectors with a unit minor. Therefore if
\[
 v_5(q^3-1)=v_5(m)+1,
\]
equation (7) forces $Z\in5L_3$.

At $q_{10}$, $m=n_{10}=14!$ has valuation $2$, while the prescribed shell has valuation $v_5(q_{10}^3-1)=3$. At $q_{15}$, $m=n_{15}=19!$ has valuation $3$ and the corresponding denominator has valuation $4$. Thus every local class-four lift would force both discrepancies in (6) to vanish. This is impossible for one value of $\lambda_E$.

Every global class-four lift would restrict to such local lifts. Hence the global class-four fibre is empty, while the integral class-three fibre was already shown nonempty. The uniqueness of the integral class-two class excludes an unexamined alternative lower fibre. This proves the theorem.
\end{proof}

\paragraph{Original certificate scope.}
The original certificate checked the complete class-three conjugation and actual $C_4$ composition formulas, the invariant covariance, the symbolic class-four tame defect, and primitive Hall minors. Its recorded exhaustive counts included $50{,}000$ actual $C_4$ compositions and $400$ invariant-covariance cases. These are historical execution records; this TeX reconstruction makes no new computational claim.
\end{document}
